% DRAFT (ojas fallback, 26 Oct): ayush to review
% =====================================================================================
% Sec. 1  Introduction.  Owner: ayush (writing lead).  Fallback draft by ojas (branch
% ojas/paper-intro) for ayush to rewrite or drop.  Budget: ~0.70 page, ~500 words, no figure.
%
% Sources: plan v2 Sec. 1 and Sec. 8 (documents/C3_Detailed_Plan_v2.pdf); Secs. 2, 3, 5 and
% limitations.tex on main; Sec. 7 on origin/ojas/paper-related (PR #24).  Since PR #34 the
% paper says nothing about pre-registration or process; keep it that way here.
% No result numbers.  The only numbers are C2 macros (artifacts/probe/macros_t9.tex) and
% design constants: k in {1,2,4,8} and rho in {0,...,1} (artifacts/probe/spec.json),
% "2--8 tokens per step" (plan v2 T10; the k >= 2 range of Sec. 5).
%
% Labels used: sec:setup, prop:split, eq:decomp (Sec. 2); sec:instrument, sec:inference (Sec. 3);
% sec:c1, sec:c2, sec:c3.  Citation keys: all in paper/refs.bib on main.
% Not mentioned: the Dorigo (2026) winner's-curse notebook (no refs.bib entry on main); if
% ayush adds it, keep it to methods (the E2 own/field split, the exact all-subsets estimator).
% =====================================================================================

\section{Introduction}
\label{sec:intro}

A masked diffusion language model writes by unmasking. Each step is one forward pass that
gives a distribution at every masked position; the decoder then commits some positions for
good. LLaDA and Dream by default commit the $k$ positions whose top token is most probable,
and threshold rules such as Fast-dLLM commit every position above a confidence
cutoff \citep{nie2025llada,ye2025dream,chang2022maskgit,wu2025fastdllm}. Confidence thus does two
jobs: it chooses which tokens are committed, and it is read as the chance that they are
right. On committed tokens it is a selected quantity. The largest of many noisy
estimates tends to be one whose error pushed it up, the winner's curse
\citep{andrews2024winners}, and a commit out of many masked positions, as at early steps and in
long outputs, is such a maximum.

The usual checks cannot tell whether this matters. Decoders are judged on accuracy and speed,
not on the calibration of the tokens they commit. Calibration, where measured, is averaged
over randomly masked positions, the states the model was trained on, not over the one the
decoder picks. Nor is the bias automatic:
selection does not bias a correct posterior \citep{efron2011tweedie}. We show that the gap
between confidence and accuracy on committed tokens splits exactly into ordinary
miscalibration, reweighted towards where winners sit, and a \emph{field term} that is nonzero
only if accuracy at fixed confidence depends on context
(Proposition~\ref{prop:split}, \S\ref{sec:setup}). Only the field term is a winner's curse,
and only an audit at the decoder's selection can see it.

We make three contributions.
\begin{itemize}
\item \textbf{C1, a selection-conditional calibration audit} (\S\ref{sec:c1}). On teacher-forced
  traces (\S\ref{sec:instrument}) from LLaDA-8B and Dream-7B, Base and Instruct, and a 110M
  masked diffusion model trained on data we control, we measure the committed-token gap as a
  function of how many candidates a token beat, its field term, and whether the states a
  decoder visits are calibrated worse than random masks. Everything fitted uses one half of
  each set's documents; every number we report comes from the held-out other half, and the
  primary contrasts are tested in Holm-corrected families across models and data
  (\S\ref{sec:inference}).
  % TODO(results): one sentence on C1, from the Holm summary (claimed stars only), keeping to
  %   the levels/trends split of limitations.tex (levels from the 110M arm, trends in m from
  %   the 8B arms).  Wording per outcome:
  %   (a) E1 and E2 claimed in some arms: "The gap grows with field size in <arms>, and part
  %       of the growth is a field term that a token's own confidence does not reveal."
  %   (b) E1 claimed, E2 not: "The gap grows with field size in <arms>, but the field term
  %       cannot be told from zero: the growth is ordinary miscalibration seen through
  %       selection, so recalibrating per step, not per context, would remove it."
  %   (c) neither claimed: "In no arm does the gap grow detectably with field size, which is
  %       why confidence thresholding works as well as it does."
  %   Add E3 (exploratory, not in a Holm family) only if its sup-t test rejects: "Decoder
  %   states are <more/less> overconfident than random masks of the same size."
\item \textbf{C2, an instrument for the two errors of parallel commitment}
  (\S\ref{sec:c2}). Committing several tokens in one step incurs a dependency error (A) and an
  estimation error (B) (Eq.~\ref{eq:decomp}). On a JSON probe with known joint distribution we
  compute both exactly, in nats, for a fine-tune of our 110M model. (A) overtakes (B) only for
  strongly coupled fields ($\rho\ge\CTwoConfCrossRhoMin$) and several tokens per step;
  ranking by confidence raises (A), not (B); and the two trade off along $k$: at $\rho=1$,
  (B) falls as (A) rises. On real text, a proxy for (A) is near zero for the 110M model and
  LLaDA-8B-Base at 2--8 tokens per step, which looks like the probe's (B)-dominated regime,
  though the proxy sees only dependence the model has learned.
\item \textbf{C3, a calibrated commitment policy} (\S\ref{sec:c3}). C3a commits every
  position whose recalibrated confidence, fitted on the fit half, implies an error rate of at
  most $\alpha$, so its threshold reads as an error rate, and it needs no extra forward pass. We
  compare it with confidence decoders at matched tokens per step in free-running generation
  with LLaDA-8B-Instruct on GSM8K and MBPP.
  % TODO(results): the C3 sentence from the T12 matched-speed table (Sec. 6).  Name the
  %   baselines actually run (final-plan.md: top-k, Fast-dLLM threshold, entropy-bounded,
  %   margin, KLASS, PC-Sampler) in place of "confidence decoders".  Wording per outcome:
  %   (a) gain: "At matched speed C3a is <more accurate than / as accurate as> the best
  %       baseline on <tasks>."
  %   (b) no gain: "At matched speed C3a is no more accurate than the Fast-dLLM threshold; we
  %       report the null."
  %   (c) T12 not done at the 30 Oct freeze: drop this item, say "We make two
  %       contributions", and move C3a to future work in limitations.tex (final-plan.md,
  %       "Cut to future work").
\end{itemize}
